% ==================== 第 6 章 前期基础与验证计划 ====================
\clearpage
\section{前期研究基础与验证计划}

\subsection{前期研究基础}

项目已形成第一阶段可运行原型。现有代码包含教师、教室、班级、课程请求和
时间槽数据模型；基于 CP-SAT 的排课求解器；独立课表校验器；候选策略与
评分服务；教师请假后的局部重排；本地 REST API；决策驾驶舱、周课表、
调课中心和资源概览四个页面。源码结构、接口样例和关键实现节选见附录。

原型采用一份小规模虚构数据进行功能验证，包含 4 位教师、3 间教室、4 个
班级、4 项课程请求和 8 个待排课次。该数据只用于验证模型和接口链路，不能
代表真实机构规模，也不据此推导行业效率或成本结论。

\subsection{当前可复核结果}

表 \ref{tab:verified-baseline} 仅列出可由当前仓库代码和自动化测试直接复核
的结果。求解耗时受设备、依赖版本和系统负载影响，因此不在此给出固定的
毫秒数；后续实验将同步记录运行环境和原始测量数据。

\begin{table}[htbp]
  \centering
  \caption{当前原型的可复核验证结果}
  \label{tab:verified-baseline}
  \small
  \begin{tabularx}{\textwidth}{L{3.0cm} L{3.2cm} Y}
    \toprule
    \textbf{验证对象} & \textbf{当前结果} & \textbf{复核依据} \\
    \midrule
    演示课次完整性 & 8 个课次均生成安排 & 求解结果课次数与课程请求课次总数一致 \\
    硬约束正确性 & 独立校验器报告 0 个冲突 & 教师、教室、班级互斥及资质、容量、设备检查 \\
    候选策略 & 学生、教师和资源效率三类方案可生成 & 方案接口及策略测试 \\
    局部重排 & 教师请假后返回新方案与差异 & 基线课表、请假约束和变更清单测试 \\
    自动化测试 & 19 项测试通过 & 12 项求解与服务测试、7 项 API 集成测试 \\
    页面与接口 & 4 个页面及主要 API 可访问 & 本地 HTTP 集成测试 \\
    \bottomrule
  \end{tabularx}
\end{table}

\subsection{实验设计}

后续实验分为五组：

\begin{enumerate}
  \item \textbf{正确性实验。}构造正常、边界和不可行实例，验证硬约束、
  数据错误提示和独立校验结果。
  \item \textbf{规模实验。}按课次数量、教师数量、教室数量和时间槽数量
  分层扩展数据，记录变量数、求解状态、首次可行解时间和总耗时。
  \item \textbf{策略敏感性实验。}逐步改变软目标权重，观察统一分项指标、
  课表结构和方案稳定性的变化，不比较不同量纲的原始目标值。
  \item \textbf{调课实验。}设置单教师请假、多时段不可用和教室停用事件，
  统计变更课次数、受影响班级和求解成功率。
  \item \textbf{可用性实验。}邀请目标用户完成数据导入、方案比较、调课和
  发布任务，记录任务完成率、操作时间和主观评价。
\end{enumerate}

\subsection{评价指标}

\begin{table}[htbp]
  \centering
  \caption{实验评价指标及计算口径}
  \label{tab:evaluation-metrics}
  \small
  \begin{tabularx}{\textwidth}{L{3.0cm} L{3.6cm} Y}
    \toprule
    \textbf{指标} & \textbf{计算口径} & \textbf{用途} \\
    \midrule
    硬冲突数 & 校验器发现的教师、教室、班级等冲突总数 & 判断方案是否可执行 \\
    课次完成率 & 已排课次除以请求课次总数 & 检查是否遗漏课程 \\
    偏好满足率 & 命中有效偏好时段的课次占比 & 衡量学生或教师体验 \\
    容量适配率 & 实际使用座位数除以分配教室容量 & 衡量资源匹配程度 \\
    负载绝对偏差 & 教师日课次相对目标值的绝对偏差之和 & 衡量日负载均衡 \\
    调课变更率 & 相对基线发生变化的课次占比 & 衡量局部重排稳定性 \\
    求解时间 & 在注明软硬件环境下的端到端耗时 & 评估规模扩展能力 \\
    任务完成率 & 用户在限定条件下完成指定操作的比例 & 评估系统可用性 \\
    \bottomrule
  \end{tabularx}
\end{table}

\subsection{结果记录与复现}

每次实验将保存输入数据、策略参数、依赖版本、随机种子、求解状态、指标和
原始输出。性能实验至少重复多次并报告中位数及波动范围；无法复现或缺少
来源的数据不进入最终结论。自动化测试在代码变更后持续运行，防止新增软目标
破坏已有硬约束。
